\documentclass{article}

% Language setting
% Replace `english' with e.g. `spanish' to change the document language
\usepackage[english]{babel}

% Set page size and margins
% Replace `letterpaper' with `a4paper' for UK/EU standard size
\usepackage[letterpaper,top=2cm,bottom=2cm,left=3cm,right=3cm,marginparwidth=1.75cm]{geometry}

% Useful packages
\usepackage{amsmath}
\usepackage{graphicx}
\usepackage[colorlinks=true, allcolors=blue]{hyperref}

\title{Software Testing 2025/6 Portfolio}
\author{s2149970}

\begin{document}
\maketitle

\section{Introduction}
This portfolio summarises the testing evidence for PizzaDronz and maps that evidence to the five learning outcomes using local documents and CI artifacts as proof of work.

\section{Outline of the Software Being Tested}
PizzaDronz is a Spring Boot REST API for autonomous pizza delivery drones. It validates orders against menu and payment rules, computes no-fly-zone-safe paths using an A* search, and returns paths (including GeoJSON). The service consumes external REST endpoints for orders, restaurants, no-fly zones, and the central area.

\section{Learning Outcomes}
\begin{enumerate}
\item \textbf{Analyze requirements to determine appropriate testing strategies} \marginpar{[default 20\%]}
\begin{enumerate}
\item \textbf{Range of requirements, functional requirements, measurable quality attributes, qualitative requirements}
Requirements are documented in \texttt{requirements.md}. Examples include functional rules (order validation: card number length, expiry format, pizza count and total; pathfinding must avoid no-fly zones and end within close distance of the destination), measurable quality attributes (path calculation completes within the 60s timeout used in tests; numeric comparisons use \texttt{EPSILON\_ERROR}), and qualitative requirements (safety: never enter no-fly zones; usability: return clear validation codes). Stakeholders include customers, operators, and regulators, which motivates both functional and safety-focused requirements.

\item \textbf{Level of requirements, system, integration, unit}
Unit testing targets geometry and validation logic (\texttt{LngLatHandlingTests}, 207 tests). Integration tests exercise pathfinding and GeoJSON creation against real no-fly zone data (\texttt{PathFindingTests} + \texttt{PathCalculationAsGeoJsonTests}, 41 tests). System-level validation uses 604 real orders pulled from the REST API (\texttt{OrderValidationTests}) and controller method tests validate endpoint logic (\texttt{ControllerTests} + \texttt{ControllerMethodTests}, 30 tests).

\item \textbf{Identifying test approach for chosen attributes}
Boundary value and equivalence partitioning are used for coordinates, angles, and card fields. Data-driven tests with \texttt{@MethodSource} check real orders and restaurants. Path safety is checked via invariants (start/end correctness, step length, no-fly avoidance, central-area constraint). Randomized \texttt{@RepeatedTest} inputs probe corner cases, and dynamic tests expand polygon coverage.

\item \textbf{Assess the appropriateness of your chosen testing approach}
The mix of boundary, data-driven, and invariant checks directly targets functional and safety requirements, and the 604-order dataset gives realistic coverage. Limitations remain: HTTP-level controller tests are absent (method tests only), external REST dependencies reduce repeatability, and concurrency, fault injection, and performance profiling are not covered.
\end{enumerate}

\item \textbf{Design and implement comprehensive test plans with instrumented code} \marginpar{[default 20\%]}
\begin{enumerate}
\item \textbf{Construction of the test plan}
\texttt{TEST\_PLAN.md} maps each requirement to a test level and suite, and places testing in a CI-based lifecycle (local run + GitHub Actions). The plan includes unit, integration, and system tests, and documents expected evidence (test reports and coverage). It also prioritises safety and payment validation first, then performance, with unit tests running per change and integration/system tests in CI.

\item \textbf{Evaluation of the quality of the test plan}
The test plan explicitly notes weaknesses: reliance on external REST endpoints, random tests without fixed seeds, and missing HTTP-level endpoint tests. These are recorded as risks and prioritised for follow-up.

\item \textbf{Instrumentation of the code}
Instrumentation includes timeouts for pathfinding via \texttt{ExecutorService} with a 60s limit, numeric tolerances using \texttt{EPSILON\_ERROR}, data scaffolding using \texttt{RestTemplate}, and CI coverage via \texttt{jacoco-maven-plugin}. Surefire XML output enables suite-level reporting.

\item \textbf{Evaluation of the instrumentation}
Instrumentation is adequate for correctness and timeout enforcement, but it does not capture per-test performance trends, memory use, or mutation adequacy. HTTP serialization and error response coverage are not instrumented.
\end{enumerate}

\item \textbf{Apply a wide variety of testing techniques and compute test coverage and yield according to a variety of criteria} \marginpar{[default 20\%]}
\begin{enumerate}
\item\textbf{Range of techniques}
Boundary value analysis, equivalence partitioning, decision tables (604 orders), invariants for path safety, combinatorial direction checks (15 of 16 compass directions), random tests (\texttt{@RepeatedTest}), negative tests (\texttt{assertThrows}), dynamic tests (\texttt{@TestFactory}), and direct controller method tests (programmatic endpoint validation).

\item\textbf{Evaluation criteria for the adequacy of the testing}
Coverage is measured with JaCoCo (instruction, branch, line) and compared to thresholds. For pathfinding, adequacy is checked by invariants (no-fly avoidance, step constraints, destination proximity). Error handling is covered by invalid card and null-input tests.

\item \textbf{Results of testing}
Latest evidence from GitHub Actions artifacts (test-results and jacoco-coverage-report in run \texttt{21045809398}): 882 tests, 0 failures. Suite times include \texttt{PathFindingTests} about 12.6s (variable), \texttt{OrderValidationTests} about 0.56s, \texttt{LngLatHandlingTests} about 0.29s, \texttt{ControllerMethodTests} about 1.7s, \texttt{ControllerTests} about 0.02s, \texttt{PathCalculationAsGeoJsonTests} about 0.50s. JaCoCo totals: instruction 90.1\%, branch 93.3\%, line 86.3\% (from the CI-generated \texttt{jacoco.xml}). Defect yield in this run is 0 (no failures), which limits confidence in untested scenarios.

\item \textbf{Evaluation of the results}
Coverage exceeds the configured thresholds (line >= 80\%, branch >= 75\%), and the order-validation and geometry suites provide strong functional confidence. Remaining risk is primarily in HTTP-level controller behavior and limited statistical sampling of paths.
\end{enumerate}

\item \textbf{Evaluate the limitations of a given testing process, using statistical methods where appropriate, and summarise outcomes} \marginpar{[default 20\%]}
\begin{enumerate}
\item \textbf{Identifying gaps and omissions in the testing process}
Key gaps: no HTTP-level endpoint tests, no load or concurrency testing, no fault injection for REST failures, random tests without fixed seeds, and no mutation testing or security checks.

\item \textbf{Identifying target coverage/performance levels for the different testing procedures}
Targets are aligned with \texttt{pom.xml}: line >= 80\% and branch >= 75\%. For performance, path calculation must complete within the 60s timeout (ideally far below). For statistical confidence in path safety, far more random start points than the current 20 are required.

\item \textbf{Discussing how the testing carried out compares with the target levels}
Line (86.3\%) and branch (93.3\%) meet targets. The 60s timeout is satisfied in all current runs, but sampling remains limited. With roughly 20 random path samples and no observed failures, the rule-of-three gives an upper bound of about 15\% failure rate at 95\% confidence, which is too weak for safety claims. HTTP-level coverage is 0\% (only method-level controller tests exist), so endpoint serialization and error responses are unverified.

\item \textbf{Discussion of what would be necessary to achieve the target level}
Add MockMvc tests for HTTP endpoints, mock external REST APIs, increase random path samples, and introduce mutation testing (PIT). Add performance baselines and concurrency tests to support stronger statistical claims.
\end{enumerate}

\item \textbf{Conduct reviews, inspections, and design and implement automated testing processes} \marginpar{[default 20\%]}
\begin{enumerate}
\item \textbf{Identify and apply review criteria to selected parts of the code and identify issues in the code}
Manual review criteria: correctness, validation, side effects, and duplication. Findings: \texttt{OrderValidationService} has a fragile expiry-date check and creates a \texttt{RestTemplate} per request; \texttt{LngLatService.validRegion} mutates the region list; \texttt{AStarPathFinding} performs REST calls inside the algorithm (harder to test and cache). These are documented as review issues rather than fixed changes.

\item \textbf{Construct an appropriate CI pipeline for the software}
GitHub Actions workflow in \texttt{.github/workflows/ci.yml} runs tests on push and PR, generates JaCoCo reports, publishes suite summaries, and uploads artifacts. JDK 18 and Maven caching are configured.

\item\textbf{Automate some aspects of the testing}
Automation includes data-driven tests, randomized and dynamic tests, CI-triggered test execution, and JaCoCo coverage checks. \texttt{@BeforeEach} in \texttt{ControllerMethodTests} constructs controllers with real services for repeatable method-level endpoint validation.

\item \textbf{Demonstrate the CI pipeline functions as expected}
CI artifacts from GitHub Actions (test results and JaCoCo coverage reports) show successful runs and coverage, and the CI workflows are configured to fail on test or coverage regressions (JaCoCo check in \texttt{pom.xml}). This provides reproducible evidence that the pipeline enforces testing standards.
\end{enumerate}
\end{enumerate}
\end{document}
